\section{Cosmic Preferred Time}
\label{crt:cpt}

The concept of a cosmic preferred time occupies a peculiar position in
modern theoretical physics. On the one hand, the principle of relativity
appears to deny the existence of any physically distinguished temporal
foliation of spacetime. On the other hand, modern cosmology naturally
introduces a cosmic time parameter associated with the large-scale
distribution of matter and radiation. Furthermore, several approaches to
quantum gravity, modified gravity, dark matter phenomenology, and
foundational quantum mechanics introduce a preferred temporal structure at
a fundamental level. This paper reviews the concept of cosmic preferred
time from historical, mathematical, and observational perspectives. We
distinguish between emergent preferred time in cosmology and fundamental
preferred time in alternative theories of gravitation. Particular attention
is given to Lorentz-violating theories, Einstein-aether models,
Hořava-Lifshitz gravity, Stueckelberg-Horwitz-Piron relativistic dynamics,
and preferred-frame interpretations of quantum theory. Current experimental
constraints and possible future tests are discussed.

Since the development of special relativity, the prevailing view has been
that no inertial frame possesses privileged physical status. The relativity
principle implies that the laws of physics are invariant under Lorentz
transformations, eliminating the need for an absolute time.

Nevertheless, cosmology appears to reintroduce a distinguished temporal
parameter. In the standard Friedmann-Lemaître-Robertson-Walker (FLRW)
universe, one defines a cosmic time $ t $  measured by observers comoving
with the Hubble flow. The remarkable isotropy of the cosmic microwave
background (CMB) identifies a unique frame in which the dipole 
anisotropy vanishes.

This raises a fundamental question:
Is cosmic time merely an emergent property of cosmological matter
distribution, or does nature possess a deeper preferred temporal structure?
The question has gained renewed importance because:
\begin{enumerate}
   \item Quantum gravity often requires a notion of evolution.
   \item Dark matter and dark energy remain unexplained.
   \item The problem of time in general relativity remains unresolved.
   \item Quantum measurement may require a preferred foliation.
\end{enumerate}

Newton assumed: $t \in \mathbb{R}$
as a universal parameter flowing uniformly throughout the universe.
Absolute simultaneity existed:$t_A=t_B$
independent of observer motion.

Lorentz and Poincaré proposed:
a preferred ether frame,
real Lorentz contraction,
hidden absolute simultaneity.
All observable predictions agreed with special relativity.
Thus:Ether theory = Special relativity ---
experimentally.

Einstein eliminated:
ether,
absolute simultaneity,
universal time.
The invariant quantity became $ds^2= g_{\mu \nu} dx^\mu dx^\nu$.
No preferred foliation exists in general.

 \subsection{Cosmic Time in Standard Cosmology}
FLRW Metric:\\
The homogeneous isotropic metric is
\begin{equation}
   \label {crt:cpt-flrw}
    ds^2=-dt^2+a^2(t) (\frac{dr^2}{1-k r^2}+r^2 d \omega^2)
\end{equation}
The parameter $ t $ is called cosmic time.
Observers at rest with respect to the Hubble flow satisfy
$ u^\mu=(1,0,0,0).$
Cosmic time corresponds to:
proper time of comoving galaxies,
age of the universe,
evolution parameter of cosmology.
Current age:$ t \approx 13.8 Gyr. $

CMB Preferred Frame\\
The cosmic microwave background defines a unique frame.
Solar system velocity:$ v \approx 369 km/s $
relative to the CMB rest frame.
This frame is physically observable.
However:
it does not violate relativity.
The laws remain Lorentz invariant.
Only the matter distribution selects the frame.

\subsection{Preferred Time in General Relativity}
A spacetime may admit a global time function $ T(x^\mu) $
such that $ \nabla_ \mu T $ is everywhere timelike.
Then hypersurfaces $ T=const $ provide a foliation.

For globally hyperbolic space-time manifolds:$ M\cong \Sigma \times \mathbb{R}.$ A global time function exists.
However:
it is not unique,
no foliation is preferred by GR itself.

Canonical GR leads to $ \hat H \Psi=0 $ (Wheeler-DeWitt equation).
No explicit time evolution appears.
This motivates attempts to introduce preferred temporal variables.

\subsection{Einstein-Aether Theory}
One of the most developed preferred-time theories is
the Einstein-aether framework.
A unit timelike vector field $ u^\mu $,
 $ u^\mu u_\mu=-1 $is introduced.
The vector field fills spacetime.
It determines a preferred local rest frame.
The action contains $ S=S_{GR}+S_{AE} $ .
Additional terms involve $ \nabla_\mu u_\nu. $
Four dimensionless coupling constants appear.

Predictions include:
modified gravitational waves,
Lorentz violation,
modified black holes,
altered cosmology.

LIGO-Virgo constraints are severe.
Observed gravitational wave speed satisfies $ |c_g -c|< 10^{-15}.$
Consequently most parameter space has been excluded.
Nevertheless small deviations remain viable.

\subsection{Hořava-Lifshitz Gravity}
Quantum gravity may require sacrificing Lorentz invariance.
Hořava proposed:$ t \rightarrow b^z t \quad x \rightarrow b x $
Time and space scale differently.
The theory assumes $ M\cong \Sigma \times \mathbb{R}.$
A preferred cosmic time exists fundamentally.
Not all coordinate transformations are allowed.
Possible benefits:
power-counting renormalizability,
improved ultraviolet behavior,
emergence of GR at large scales.
Problems include:
extra scalar graviton,
strong coupling,
observational constraints.
The ultimate viability remains uncertain.

\subsection{ Preferred Time and Dark Matter}
An intriguing possibility is that dark matter phenomena arise from fields associated with preferred temporal structures.
Several proposals suggest:
vector-field dark matter,
aether condensates,
superfluid dark matter,
emergent gravity scenarios.
In Einstein-aether cosmology, perturbations of the preferred frame can mimic some dark-sector effects.

However no fully successful replacement for dark matter currently exists.
Evidence from:
galaxy clusters,
lensing,
CMB anisotropies,
still strongly favors genuine dark matter.

\subsection{Preferred Time in Quantum Foundations}
Relativistic Bohmian theories often require $ \Sigma_t $
hypersurfaces defining simultaneity.
A preferred foliation naturally appears.

Theories such as spontaneous collapse frequently introduce a universal temporal ordering.
Collapse occurs on preferred hypersurfaces.

In the SHP framework, evolution occurs in an invariant parameter
$ \tau $. The parameter is independent of coordinate time:$ x^\mu(\tau). $
This introduces a form of universal dynamical ordering.
The SHP parameter is not cosmic time itself, but it plays a role analogous
to a universal evolution parameter.
One may interpret:$ \tau $
as a generalized preferred time.
This is one of the strongest motivations for investigating cosmic temporal
structures beyond GR.

\subsection{ Experimental Searches}
\begin{description}
    \item [Michelson-Morley Type Tests ] Modern optical resonators constrain anisotropy:$ \Delta c/c<10^{-18} .$ No preferred frame effects observed.
    \item [Atomic Clocks ] Comparisons of ultra-precise clocks constrain Lorentz violation.Current bounds are extraordinarily strong.
    \item [Cosmic Rays ] Ultra-high-energy cosmic rays limit modified dispersion relations.
    \item [Gravitational Waves ] Multi-messenger observations provide the
        strongest modern constraints. The event $GW170817$ showed
        gravitational and electromagnetic signals
arriving almost simultaneously.
Many preferred-frame theories were ruled out.
\end{description}
\subsection{ Relation to Global Time Programs}
Several independent research directions seek a universal temporal structure:
\begin{description}
    \item [ York Time] $ T_Y \propto K $ (trace of extrinsic curvature).
    \item [ Shape Dynamics] Replaces spacetime symmetry by conformal symmetry. 
        A preferred slicing emerges naturally.
    \item [ CMC Foliations] Constant-mean-curvature hypersurfaces provide distinguished time coordinates.
    \item [ Thermal Time] Proposed by Carlo Rovelli and collaborators.
Time emerges from statistical states rather than fundamental geometry.
\end{description}
\subsection{ Philosophical Interpretation}
Three possibilities exist:
\begin{description}
    \item [ View A:]  No Preferred Time
Fundamental reality is fully four-dimensional.
All preferred times are gauge artifacts.
This is the orthodox GR position.
    \item [ View B:]  Emergent Preferred Time
Cosmic time exists because matter selects a frame.
The CMB frame is real but not fundamental.
This is currently the mainstream cosmological view.
    \item [ View C:]  Fundamental Preferred Time
A deeper temporal structure exists.
Relativity emerges approximately.
Examples include:
Einstein-aether theory,
Hořava gravity,
SHP dynamics,
certain collapse theories.
This remains speculative but experimentally testable.
\end{description}

